% ============================== Part F ==============================
\problem{F1}{The Laplace approximation}

\begin{given}
Near its maximum $\mathbf x_0$, a log density is approximately quadratic,
$\ln P(\mathbf x) \simeq \ln P(\mathbf x_0) - \tfrac12(\mathbf x - \mathbf x_0)^\T\mathsf A(\mathbf x - \mathbf x_0)$, with
$\mathsf A = -\nabla\nabla\ln P|_{\mathbf x_0}$. Treating $P$ as a Gaussian with covariance $\mathsf A^{-1}$ is the
Laplace approximation.
\source{\cMacKayLaplace.}
\end{given}

\need{A1 (a Gaussian's log is a quadratic).}

\begin{enumerate}
  \item For $P(x) \propto x^a e^{-bx}$ ($x > 0$), find the mode, $\mathsf A$, and the Laplace standard deviation.
  \item The exact standard deviation is $\sqrt{a+1}/b$. Compare for $a = 4$. When does Laplace work well?
  \item For the posterior of E2, show that the Laplace approximation is exact, with $\mathsf A = \mathsf H$.
\end{enumerate}

\solution{
(a) $\ln P = a\ln x - bx$; the mode is $x_0 = a/b$, $\mathsf A = a/x_0^2 = b^2/a$, so $\sigma_L = \sqrt a/b$.
(b) $a = 4$: $2/b$ against $\sqrt5/b = 2.24/b$, 11\% low. It works well when the distribution is close to Gaussian,
here for large $a$ (the ratio $\sqrt{a/(a+1)} \to 1$).
(c) The log posterior is exactly quadratic, with $-\nabla\nabla\ln p = S^{-1} + A^\T N^{-1}A = \mathsf H$.
}

% --------------------------------------------------------------------
\problem{F2}{Gauss--Newton: curvature from the model, not from the data grid}

\begin{given}
For a least-squares objective $F(\mathbf x) = \tfrac12\sum_i f_i(\mathbf x)^2$, the exact Hessian is
$J^\T J + \sum_i f_i\,\nabla\nabla f_i$, with $J$ the Jacobian of the residuals $f_i$. Gauss--Newton keeps only $J^\T J$.
\source{\cGaussNewton.}
Our galaxies enter through $\ln L_i(\vg_i(\kappa))$, read off a grid by bilinear interpolation, and each $L_i$ has
precision $F_i$ (A3).
\end{given}

\need{A3, F1.}

\begin{enumerate}
  \item Suppose $\ln L_i(\vg) \approx -\tfrac12(\vg - \hat\vg_i)^\T F_i(\vg - \hat\vg_i)$. Show that the Hessian of
    $-\sum_i\ln L_i(\vg_i(\kappa))$ with respect to the pixels is
    $\sum_i J_i^\T F_iJ_i + (\text{terms} \propto \text{residuals} \times \nabla\nabla\vg_i)$, with
    $J_i = \partial\vg_i/\partial\kappa$.
  \item Why are the dropped terms small for us? (Think about how $\vg_i$ depends on $\kappa$; see Part D.)
  \item A bilinear interpolant is linear along each axis inside a grid cell. What is its second derivative along $g_1$
    there? Why does that make the \emph{exact} Hessian of our objective useless, and why does $J^\T FJ$ not have this
    problem?
\end{enumerate}

\solution{
(a) The gradient is $-\sum_i J_i^\T\nabla_{\vg}\ln L_i = \sum_i J_i^\T F_i(\vg_i - \hat\vg_i)$. Differentiating again gives
$\sum_i J_i^\T F_iJ_i$ plus terms with $\partial J_i/\partial\kappa$ multiplied by the residual $F_i(\vg_i - \hat\vg_i)$.
This is the weighted version of MNT04's Eq.~(3.4b).
(b) $\vg_i$ is linear in $\kappa$ except for the $1/(1-\kappa)$ factor (Part D), so $\nabla\nabla\vg_i$ is small, and it is
multiplied by residuals that scatter around zero.
(c) Zero. Inside each grid cell the interpolant has no curvature along $g_1$ or $g_2$; the curvature lives only in the
kinks between cells. The exact Hessian therefore misses most of the data's information and can even have negative
directions (in our first attempt, 12 pixels got $\sigma = 0$). $J^\T FJ$ takes the curvature from each galaxy's known
precision $F_i$, which is positive definite by construction.
}

% --------------------------------------------------------------------
\problem{F3}{Two different uncertainties, and the resolution matrix}

\begin{given}
In the linear-Gaussian model of E2, write $\mathsf H_d = A^\T N^{-1}A$ (data) and $\mathsf H_p = S^{-1}$ (prior), so
$\mathsf H = \mathsf H_d + \mathsf H_p$ and $\hat{\mathbf s} = \mathsf H^{-1}A^\T N^{-1}\vd$. Let the truth be
$\mathbf s_{\rm true}$, with $\vd = A\mathbf s_{\rm true} + \mathbf n$. \emph{Derived here}
(\texttt{KAPPA\_FITTING\_EXPLAINED.md}, Sec.~8.3).
\end{given}

\need{E2, F1.}

\begin{enumerate}
  \item Show that $\E_{\mathbf n}[\hat{\mathbf s}] = \mathsf R\,\mathbf s_{\rm true}$ with $\mathsf R = \mathsf H^{-1}\mathsf H_d$.
  \item Show that the covariance of $\hat{\mathbf s}$ over noise realizations is $\mathsf H^{-1}\mathsf H_d\mathsf H^{-1}$.
  \item Scalar case ($\mathsf H_d = h$, $\mathsf H_p = p$): compare this noise variance with the posterior variance
    $1/(h+p)$. Which is larger, and what does each one answer?
  \item On our map, the noise covariance predicts an aperture scatter of 0.0034 and the posterior covariance 0.0067. In
    the scalar picture, what data-to-prior ratio $h/p$ does that imply? For ``is this blob a noise fluctuation?'',
    which one should the S/N map use?
\end{enumerate}

\solution{
(a) $\E[\vd] = A\mathbf s_{\rm true}$, so \[\E[\hat{\mathbf s}] = \mathsf H^{-1}A^\T N^{-1}A\,\mathbf s_{\rm true} = \mathsf H^{-1}\mathsf H_d\,\mathbf s_{\rm true}.\]
(b) $\mathrm{Cov} = \mathsf H^{-1}A^\T N^{-1}\,N\,N^{-1}A\mathsf H^{-1} = \mathsf H^{-1}\mathsf H_d\mathsf H^{-1}$.
(c) Noise variance $h/(h+p)^2$ = posterior variance $\times\, h/(h+p)$, so it is always smaller. The posterior variance
answers ``how unsure am I about the true value?'', including what the prior allows. The noise variance answers ``how
much would this map change with new noise?''.
(d) $(0.0034/0.0067)^2 = 0.26 = h/(h+p)$, so $h/p \approx 0.35$: on the subhalo's scale the prior outweighs the data
about 3 to 1. The S/N map should use the noise covariance, which is what checks out against the scatter over
synthetic realizations (0.0029).
}

% --------------------------------------------------------------------
\problem{F4}{\py\ How many numbers do the data actually determine?}

\begin{given}
In the eigenbasis of the data Hessian (eigenvalues $\lambda_a$), with an isotropic prior of precision $\alpha$, the number
of parameters well determined by the data is $\gamma = \sum_a \lambda_a/(\lambda_a + \alpha)$. Each term, between 0 and
1, measures the data's strength relative to the prior in that direction.
\source{\cMacKayGamma.}
\end{given}

\need{F3 (the resolution matrix $\mathsf R$), E5 (the 1D toy problem).}

\begin{enumerate}
  \item Show that $\mathrm{tr}\,\mathsf R = \gamma$ when $\mathsf H_p = \alpha I$. (Diagonalise $\mathsf H_d$.)
  \item \py\ For the 1D problem of E5, compute $\mathrm{tr}\,\mathsf R$ as a function of the number of pixels (keeping
    the physical $\ell$ fixed). Does it keep growing?
  \item Our tuned map has 144 pixels and about 150 galaxies, and $\mathrm{tr}\,\mathsf R = 16.5$. With a weaker prior
    ($\sk = 0.1$) it is 53; with 576 pixels it is 10. Answer the question ``how can 141 galaxies fit 144 pixels?''
  \item What sets the map's effective resolution, if not the pixel size?
\end{enumerate}

\solution{
(a) With $\mathsf H_d = U\Lambda U^\T$, $\mathsf R = (\mathsf H_d + \alpha I)^{-1}\mathsf H_d = U\,\mathrm{diag}\big(\lambda_a/(\lambda_a+\alpha)\big)U^\T$,
whose trace is $\gamma$.
(b) No. In the notebook $\mathrm{tr}\,\mathsf R$ rises to $\approx 19$ near 60 pixels and then \emph{falls} (7.7 at
240). Extra pixels add only directions the data cannot see. And with $\sk$ fixed per pixel, the amplitude term
$\sum\kappa_j^2$ runs over more pixels, so every smooth pattern is penalised more: a finer grid with the same $\sk$ is a
stronger prior.
(c) They don't, and they don't need to. It is a regularised fit: the galaxies pin down about 16 independent patterns,
and the prior fixes the remaining $\sim128$. A weaker prior lets the data set more patterns, but noisily. More pixels
add only prior-dominated directions (and, at fixed $\sk$, a stronger prior), as the drop to 10 at 576 pixels shows.
(d) The rows of $\mathsf R$: their width is the resolution and their height the shrinkage, both set jointly by the
galaxies (number, positions, precisions) and by the prior. The pixel size only sets the finest scale the map could
show.
}
